\bookchapter{第三十三章：未來展望——手機的下一程}

\chapterintro{手機不會一夜之間消失，但它的“獨霸地位”可能被新一代裝置分流。這一章我們聊三件事：下一代裝置長什麼樣、手機如何“隱身”為新身份基石、以及那些我們必須提前思考的倫理議題。}

\booksection{一、AR/VR與空間計算：手機之外的“第二桌面”}

下一代裝置的競爭已經打響，最有可能分流手機時間的，是AR/VR與空間計算裝置：

\begin{itemize}
  \item \textbf{Vision Pro後續}：蘋果Vision Pro的釋出標誌著"空間計算"概念的正式落地，但第一代產品價格（29999元起）和重量仍是門檻。預計2026-2027年會推出更輕、更便宜的平價版Vision，讓空間計算從"開發者玩具"走向"早期大眾"。屆時iPhone可能成為Vision眼鏡的算力伴侶與資料中樞。
  \item \textbf{AR眼鏡}：給你的眼前開個"第二桌面"，數字資訊變得像空氣一樣無處不在。Meta與Ray-Ban合作的智慧眼鏡已賣出數百萬副，國內Rokid、XREAL、INMO等廠商也在推出消費級AR眼鏡。2025-2026年的AR眼鏡已能做實時翻譯、導航HUD、AI助手對話，但仍受限於光波導顯示和續航，距離"取代手機"還有3-5年。
  \item \textbf{智慧戒指與可穿戴}：Samsung Galaxy Ring、Oura Ring等智慧戒指開始流行，作為手機的"貼身感測器"補充健康監測。手機的角色正從"萬能終端"轉向"中樞+多個輕量化配件"的組合。
\end{itemize}

\booksection{二、腦機介面：思想成為互動媒介}

思想本身將成為互動的媒介，但這條路比想象中漫長：

\begin{itemize}
  \item \textbf{Neuralink最新進展}：馬斯克的Neuralink已於2024年完成首例人體植入，首位受試者Noland Arbaugh透過意念控制電腦滑鼠下國際象棋。目前植入式BCI（腦機介面）主要面向癱瘓患者輔助場景，目標是幫助他們重新"操作"數字世界。
  \item \textbf{非侵入式BCI}：多家公司開發非侵入式頭環（如Kernel、OpenBCI），無需手術即可讀取腦電訊號，但精度遠低於侵入式。距離"用意念打字"還有相當距離。
  \item \textbf{消費級展望}：腦機介面距離消費級手機互動還有10-20年，但短期可能先用於無障礙輔助、專注力監測等場景。倫理與隱私問題（"我的腦電波算誰的隱私"）將先於技術成熟成為爭論焦點。
\end{itemize}

\booksection{三、手機的“隱身”與新身份基石}

手機的另一個未來進化方向或許是“隱身”——化身為無處不在的“幕後推手”，成為你數字身份的唯一基石：

\begin{itemize}
  \item \textbf{雲端雙因素認證}：瞬間完成"你是你"的認證，登入任何服務無需手動輸驗證碼。
  \item \textbf{數字遺產成為出廠設定}：新機啟用時預設數字資產繼承方案，讓"身後事"成為出廠預設。
  \item \textbf{從"裝置中心"到"身份中心"}：你的雲端個人資料保險櫃才是中心，硬體只是訪問入口。換機、換裝置，身份與資料始終跟著你。
\end{itemize}

\booksection{四、倫理思考：我們，還是手機的主人嗎？}

一個終極拷問始終存在：我們是被技術賦能的“神”，還是被演演算法餵養的“缸中之腦”？

\begin{itemize}
  \item \textbf{AI倫理}：當端側AI比你自己更懂你的偏好，你的"選擇"是自主意志還是被引導的結果？AI Agent代你下單、回覆訊息時，責任與邊界在哪裡？
  \item \textbf{資料主權}：你的健康資料、腦電資料、生物特徵資料，是屬於你、屬於裝置廠商還是屬於雲服務商？《個人資訊保護法》和歐盟GDPR只是起點，"資料可攜帶權"和"被遺忘權"的真正落地還有很長的路。
  \item \textbf{數字鴻溝}：AI手機、摺疊屏、衛星通訊都在拉大"用得起"和"用不起"之間的差距。數字基礎設施的公平接入，將成為新的社會議題——就像20年前"會不會上網"一樣關鍵。
  \item 當演演算法比你更懂你想看什麼，你的"喜歡"有多少是出於自發？
  \item 當你每隔幾分鐘就不自覺地點亮螢幕，這是你的習慣還是被設計出來的"上癮行為"？
\end{itemize}

\booksection{五、保持連線，也保持獨立——你的數字生活哲學}

無論技術如何進化，下面這幾條數字生活哲學，10年後依然適用：

\begin{enumerate}
  \item 做資源的分配者，而非被推送的接收者。
  \item 讓自動化成為你忠實的"副駕駛"。
  \item 建立你穩固的"數字後方"。
  \item 定期執行"數字斷舍離"與"社交齋戒"。
  \item 永遠把安全和隱私作為一切便利的前提。
\end{enumerate}

\begin{tcbpractice}{來試試：制定你的《個人數字裝置使用服務條款》}
為你和你的手機（以及未來的AR眼鏡、AI助手）寫一份簡單的“服務條款”，來定義你們的關係。例如：

\begin{itemize}
  \item 第1條：我的注意力時間屬於我。你在我工作/陪伴家人時應主動進入"勿擾"狀態。
  \item 第2條：我的個人資料是我的財產。你的首要職責是安全加密與便捷同步。
  \item 第3條：我們的共同目標是提升生活效率與體驗。若你透過誘導性設計讓我沉迷，我有權將你解除安裝。
\end{itemize}

\end{tcbpractice}

\begin{tcbsummary}{本章小結}
\begin{itemize}
  \item Vision Pro與AR眼鏡是手機之外的新計算入口，但短期難以取代手機
  \item 腦機介面距離消費級還有10-20年，倫理問題先於技術成熟
  \item 手機正向"隱身"方向進化，從裝置中心走向身份中心
  \item AI倫理、資料主權、數字鴻溝是未來10年的核心議題
  \item 保持連線也保持獨立，是數字生活不變的哲學
\end{itemize}

\end{tcbsummary}
